\chapter{Elliptische Kurven über endlichen Körpern}
Für kryptographische Verfahren betrachtet man elliptische Kurven über endlichen Körpern. Dort lassen sich die Körperoperationen exakt ausführen, und die Kurvenpunkte bilden eine endliche Gruppe (vgl. \cite{omondi20}, S.~230). In diesem Kapitel werden zunächst endliche Körper und die Punktgruppe elliptischer Kurven untersucht. Anschließend werden Eigenschaften und Verfahren behandelt, die für die Auswahl geeigneter Kurven wichtig sind. Wo die Charakteristik des Körpers weder $2$ noch $3$ ist, wird die bereits eingeführte kurze Weierstraß-Gleichung verwendet.
\section{Endliche Körper}
In diesem Abschnitt werden die Grundlagen zu endlichen Körpern gelegt.
Die Darstellung orientiert sich an Stroth und Waldecker
\cite[S.~66--67, 163--164]{stroth24}, Karpfinger
\cite[S.~46, 317--318, 391--392]{karpf24} sowie Güneysu,
Paar und Pelzl \cite[S.~279~f.]{guneysu24}.
\begin{definition}
    Ein endlicher Körper ist ein Körper mit endlich vielen Elementen. Für jede Primzahlpotenz \(q=p^n\) existiert bis auf Isomorphie genau ein endlicher Körper mit \(q\) Elementen. Dieser wird dann mit \(\mathbb F_q\) oder \(GF(q)\)(Galois Field) bezeichnet.
\end{definition}
In kryptographischen Anwendungen treten insbesondere die Körper \(\mathbb F_p\) und \(\mathbb F_{2^m}\) auf. Der kleinste Durchschnitt $P$ aller Teilkörper von $K$ wird der \textbf{Primkörper} genannt
\[
    P=\cap \{U|U \text{ist ein Teilkörper von }K\}.
\]
\begin{satz}

    \label{hilfs:primkörper}
    Seien $K$ ein Körper und $K_0$ der Primkörper von $K$.
    \begin{enumerate}
        \item Ist char($K)=0$, so ist $K_0$ zu $\mathbb{Q}$ isomorph.
        \item Ist char($K)=p\neq 0$, so ist $K_0$ zu $\mathbb{Z}/p\cdot \mathbb{Z}$ isomorph.
    \end{enumerate}
\end{satz}
Mit Worten heißt dies: Hat $K$ die Charakteristik 0, so ist der Primkörper von $K$ zu $\mathbb{Q}$ isomorph. Ist die Charakteristik $p \neq 0$, so ist er zu $\mathbb{Z}/p$ isomorph.
\begin{proof}
    \begin{enumerate}
        \item Sei \(\Phi:\mathbb{N}\leftarrow K_0\) eine Abbildung für alle $n\in\mathbb{N}$. Sei
              \[
                  n^\Phi = (1_k+\cdots+1_k). \quad \text{n-mal}.
              \]
              Sind $n,m\in\mathbb{N}$ und $n\leq m$ mit der Eigenschaft $n^\psi=m^\psi$, so folgt
              \[
                  0_K=n^\psi-m^\psi=(1_K+\cdots+1_K). \quad \text{(n-m)-mal}.
              \]
              Da char($K)=0$. ist ist, dann $n=m$. Also ist $\psi$ eine injektive Abbildung. Somit besitzt jede $K_0$ eine Teilmenge, die sich bezüglich der Addition wie $\mathbb{N}$ verhält. Da $K_0$ ein Körper ist, gibt es also in $K_0$ einen zu $\mathbb{Z}$ isomorphen Teilring und dann auch einen zu $\mathbb{Q}$ isomorphen Teilkörper. Da $\mathbb{Q}$ keine echten Teilkörper hat, ist $K_0$ selbst isomorph zu $\mathbb{Q}$.
        \item Seien nun $p$ eine Primzahl und char($K)= p$. Sei weiter
              \[
                  L = {0_K,1_K,\dots,1_k+\cdots+1_K}\quad \text{(p-1)-mal}
              \]
              Da char($K)=p$ ist, sind die Elemente in $L$ paarweise verschieden und daher ist $|L|=p$, außerdem ist $L$ unter Addition abgeschlossen.

    \end{enumerate}
\end{proof}

\begin{satz}
    Sei $K$ ein endlicher Körper. Dann gibt es eine Primzahl $p$ und ein $n\in \mathbb{N}$, sodass \(|K|=p^n\) ist.
\end{satz}

\begin{proof}
    Da \(K\) endlich ist, müssen in der Folge
    \[
        0,\ 1_K,\ 1_K+1_K,\ 1_K+1_K+1_K,\ \dots
    \]
    zwei Elemente übereinstimmen. Daher gibt es eine kleinste positive
    natürliche Zahl \(p\) mit
    \[
        p\cdot 1_K=0.
    \]
    Diese Zahl ist die Charakteristik von \(K\).

    Die Zahl \(p\) ist eine Primzahl. Wäre nämlich \(p=rs\) mit
    \(1<r,s<p\), so wäre
    \[
        (r\cdot 1_K)(s\cdot 1_K)=p\cdot 1_K=0.
    \]
    Da ein Körper keine Nullteiler besitzt, müsste
    \(r\cdot 1_K=0\) oder \(s\cdot 1_K=0\) gelten. Dies widerspricht
    der Minimalität von \(p\).

    Der von \(1_K\) erzeugte Primkörper besitzt somit \(p\) Elemente
    und ist isomorph zu \(\mathbb{F}_p\). Folglich ist \(K\) ein
    Vektorraum über \(\mathbb{F}_p\). Da \(K\) endlich ist, besitzt
    dieser Vektorraum eine endliche Dimension \(n\).

    Bezüglich einer Basis besitzt jedes Element von \(K\) eine
    eindeutige Darstellung
    \[
        a_1v_1+\dots+a_nv_n
        \qquad\text{mit }a_i\in\mathbb{F}_p.
    \]
    Für jeden der \(n\) Koeffizienten gibt es \(p\) Möglichkeiten.
    Daher gilt
    \[
        |K|=p^n.
    \]
\end{proof}


\begin{lemma}
    \label{lem:a^q=a}
    Für einen endlichen Körper K der Charakteristik p mit $q=p^n$ gilt
    \[
        a^q=a
    \] für jedes $a\in K$.
\end{lemma}

\begin{proof}
    Für \(a=0\) ist die Aussage unmittelbar erfüllt, da $0^{p^n}=0$. Sei daher \(a\in K^\times =K\setminus\{0\}\).
    Wegen $|K^\times|=p^n-1$ gilt nach dem kleinen Satz von Fermat vgl. \cite{karpf24} S.46 gilt  in einer endlichen Gruppe $G$, \(a^{|G|} = e\) für alle \(a\in G\). Dann \(a^{p^n-1}=1\) für jedes $a\in K^\times$. Daraus folgt dann $a^{p^n}=a$ für alle $a\in K $.
\end{proof}
Diese Aussage kann nun auch noch in beide Richtungen formuliert werden, da wir sie später bei einem Beweis benötigen.
\begin{lemma}
    Für \(a\in \overline{\mathbb{F}}_q\) gilt
    \[
        a^q=a\Leftrightarrow a\in\mathbb{F}_q
    \]
\end{lemma}
\begin{proof}
    Nach dem vorherigen Lemma gilt für jedes
    \(a\in\mathbb F_q\), dass \(a^q=a\), also
    \(a^q-a=0\). Somit sind die \(q\) Elemente von
    \(\mathbb F_q\) verschiedene Nullstellen des
    Polynoms \(f(X)=X^q-X\).

    Da \(f\) den Grad \(q\) besitzt, hat es auch über
    dem algebraischen Abschluss \(\overline{\mathbb F}_q\)
    höchstens \(q\) verschiedene Nullstellen. Seine
    Nullstellen sind daher genau die Elemente von
    \(\mathbb F_q\). Folglich gilt für jedes
    \(a\in\overline{\mathbb F}_q\):
    \[
        a^q=a \quad\Longleftrightarrow\quad a\in\mathbb F_q.
    \]
\end{proof}
\begin{lemma}
    \label{lem:freshmans-dream}
    Seien $p$ eine Primzahl, $n\in\mathbb{N}$ und $q=p^n$. Sei Körper $K$ der Charakteristik $p>0$. Sind $a,b\in K$, so  ist
    \[
        (a+b)^q=a^q+b^q.
    \]
\end{lemma}
\begin{proof}
    Seien $p$ eine Primzahl, $n\geq 1$ und $q=p^n$. Für alle $a,b\in K$ gilt nach dem binomischen Lehrsatz gilt
    \[
        (a+b)^p=\sum_{i=0}^p\binom{p}{i}a^{p-i}b^i= a^p+\binom{p}{1}a^{p-1}b+\cdots+\binom{p}{p-1}ab^{p-1}+b^p
    \]
    Für $1\leq i\leq p-1$ ist
    \[
        \binom pi=\frac{p!}{i!(p-i)!}
    \]
    durch $p$ teilbar. Denn der Zähler enthält den Faktor $p$, wobei der Nenner diesen nicht enthält. Da $K$ Charakteristik $p$ hat,
    verschwinden somit alle mittleren Summanden. Es folgt
    \[
        (a+b)^p=a^p+b^p.
    \]
    Wendet man diese Gleichung $n$-mal an, erhält man
    $(a+b)^{p^n}=a^{p^n}+b^{p^n}$
\end{proof}

\begin{definition}[Elliptische Kurve über \(\mathbb F_q\)]
    Sei \(p>3\) eine Primzahl mit $q=p^r$. Eine elliptische Kurve über \(\mathbb F_q\) ist die Menge von allen Punkten $(x,y)\in \mathbb{F}_q$ mit
    \[
        y^2 \equiv x^3+Ax+B \pmod p
    \]
    und einem Punkt im Unendlichen.
    \[
        4a^3+27b^2 \not\equiv 0 \pmod p
    \] gilt. Die Bedingung stellt wie im allgemeinen Fall sicher, dass die Kurve nichtsingulär ist. Die Einschränkung $p>3$ erlaubt die Verwendung der kurzen Weierstraß-Gleichung. Für Körper der Charakteristik \(2\) und \(3\) haben wir bereits in Kapitel 3 gesehen, dass die allgemeine Weierstraß-Gleichung verwendet werrden muss.
\end{definition}

\section{Die Punktmenge $E(\mathbb{F}_q)$}
In diesem Kapitel wird der Frage nachgeganngen wie die Gruppe $E(\mathbb{F}_q)$ entsteht und wie man ihre Größe bestimmen kann. Die Darstellung dieses Abschnitts orientiert sich an Washington \cite[S.~77--79, 95 und 104--105]{washington08}, Silverman
\cite[S.~29]{silverman06}, Müller-Stach und Piontkowski
\cite[S.~51~ff.]{muller} sowie Güneysu, Paar und Pelzl
\cite[S.~285~f.]{guneysu24}.
\subsection{Explizite Punktbestimmung}
Sei E die Kurve $y^2=x^3+x+1$ über $\mathbb{F}_5$. Um nun die Punkte von E zu erhalten, berchnen wir zuerst $x^3+x+1\mod 5$ und anschließend die Quadratwurzel von diesem Polynom über $F_5$.
\begin{center}
    \begin{tabular}{c|c|c|c}
        x & $x^3+x+1$ & y       & Points       \\
        \hline
        0 & 1         & $\pm$ 1 & (0,1), (0,4) \\
        1 & 3         & -       & -            \\
        2 & 1         & $\pm$ 1 & (2,1),(2,4)  \\
        3 & 1         & $\pm 1$ & (3,1), (3,4) \\
        4 & 4         & $\pm 2$ & (4,2),(4,3)
    \end{tabular}
\end{center}
Es gilt $-1\equiv 4 \mod 5$ und $ -2\equiv 3 \mod 5$ und für $y^2=3$ gibt es keine Wurzel in $F_5$. Also lässt sich die Ordnung der Gruppe durch Punktzählung bestimmen,
\[
    \mid E(\mathbb{F}_5)\mid = \# \{(x,y)\in \mathbb{F}_5 \mid y^2=x^3+x+1\}+1 = 9
\]
Nun soll die Addition von zwei Punkte $P_1=(3,1)$ und $P_2=(2,4)$ mit Hilfe der Gruppengesetze aus Kapitel 3 auf $E(\mathbb{F}_q)$ durchgeführt werden
\[
    \frac{4-1}{2-3}=\frac{3}{-1}=\frac{3}{4}=3\cdot4^{-1} =3\cdot 4=12  \equiv 2 \mod 5.
\]
Es ist also $y=2(x-3)+1\equiv 2x$. Wenn wir dass nun in die Gleichung für die Kurve einsetzen erhalten wir
\[
    (2x)^2=x^3+x+1\Leftrightarrow 0 = x^3-4x^2+x+1.
\]
Wir kennen schon zwei Schnittpunkte der Kurve $x=3$ und $x=2$, da $(3,1)$ und $(2,4)$ auf der Gerade und Kurve liegen. Da das Polynom kubisch ist, muss die Summe der drei Nullstellen 4 sein. Damit gilt,
\[
    3+2+r \equiv 4 \mod 5 \Leftrightarrow r\equiv 4 \mod 5
\]
Also ist die dritte Koordinate $x=4$, nun erhalten wir $y=2x\equiv 3$, gespiegelt über die x-Achse erhält man die Summe
\[(3,1)+(2,4)=(4,2)\]


\subsection{Quadratische Reste und Legendre-Symbol}
Für die Punktbestimmung muss die Gleichung $y^2\equiv a\mod p$ oft gelöst werden und somit ein quadratischer Rest Modulo bestimmt werden. Bei kleinen Primzahlen ist das noch durch überlegen und ausprobieren recht einfach möglich, doch bei größeren Resten wird das deutlich problematischer. Es gibt aber ein praktisches Werkzeug um zu bestimmen ob eine Zahl ein quadratischer Rest ist oder nicht, das Legendre-Symbol.
\begin{definition}[Legendre-Symbol]
    Sei \(p\) eine ungerade Primzahl und \(a\in\mathbb Z\). Das Legendre-Symbol ist definiert durch

    \[
        \left(\frac{a}{p}\right)
        =
        \begin{cases}
            1,  & \text{falls } a \not\equiv 0 \pmod p \text{ ein quadratischer Rest modulo } p \text{ ist}, \\
            -1, & \text{falls } a \text{ ein quadratischer Nichtrest modulo } p \text{ ist},                 \\
            0,  & \text{falls } a \equiv 0 \pmod p.
        \end{cases}
    \]
\end{definition}
Es gibt zwei Methoden um das Legendre-Symbol effizient zu berechnen:
\begin{itemize}
    \item Das Eulersche Kriterium
    \item Das Quadratische Reziprozitätsgesetz
\end{itemize}
Auf beide wird nicht genauer eingeganngen, aber sie sind wichtig um Punkte für jede beliebige elliptische Kurve zu bestimmmen, da dadurch überprüft werden kann, ob es sich um einen quadratischen Rest oder nicht handelt.
Denn \(y^2=x^3+Ax+B\)
In \cite{muller} S.51ff. Kapitel 8 kann mehr zu dem Thema quadratische Reste nachgelesen werden und effiziente Methoden sowie die Beweise dazu gefunden werden.
Das Legendre-Symbol kann verallgemeinert werden über alle endlichen Köprer $\mathbb{F}_q$.
\begin{definition}
    Sei \(\mathbb F_q\) ein endlicher Körper mit \(q\) ungerade. Für \(x\in\mathbb F_q\) definieren wir
    \[
        \left(\frac{x}{\mathbb F_q}\right)
        =
        \begin{cases}
            1,  & \text{falls } t^2=x \text{ eine Lösung } t\in\mathbb F_q^\times \text{ besitzt}, \\
            -1, & \text{falls } t^2=x \text{ keine Lösung } t\in\mathbb F_q \text{ besitzt},       \\
            0,  & \text{falls } x=0.
        \end{cases}
    \]

\end{definition}
Diese Definition kann dann genutzt werden um eine erste Berechnung der Punktmenge zu ermöglichen.
\begin{theorem}
    Sei $E$ eine elliptische Kurve definiert mit \(y^2=x^3+Ax+B\) über $\mathbb{F}_q$. Dann gilt
    \[
        \#E(\mathbb{F}_q)=q+1+\sum_{x\in\mathbb{F}_q}\left(\frac{x^3+Ax+B}{\mathbb{F}_q}\right)
    \]
\end{theorem}
\begin{proof}
    Sei $x\in\mathbb F_q$ fest und sei $a\in\mathbb{F}_q$ eine Lösung der Weierstraßgleichung $a=x^3+Ax+B$. Da $q$ ungerade ist, besitzt $y^2=a$ genau zwei Lösungen in $\mathbb F_q$, wenn $a$ ein von null verschiedenes
    Quadrat ist, genau eine Lösung, wenn $a=0$ ist, und keine Lösung, wenn $a$ kein Quadrat ist. Die Anzahl der Punkte mit erster Koordinate $x$ beträgt daher
    \[
        1+\left(\frac{x^3+Ax+B}{\mathbb F_q}\right).
    \]
    Durch Summieren über alle $x\in\mathbb F_q$ und Hinzufügen des Punktes im Unendlichen $\mathcal O$ folgt
    \[
        \#E(\mathbb F_q)
        =1+\sum_{x\in\mathbb F_q}
        \left(1+\left(\frac{x^3+Ax+B}{\mathbb F_q}\right)\right)
        =q+1+\sum_{x\in\mathbb F_q}
        \left(\frac{x^3+Ax+B}{\mathbb F_q}\right).
    \]
\end{proof}
\begin{beispiel}
    Es sei $E$ eine elliptische Kurve über dem Köprer $\mathbb{F}_{11}$ mit  $y^2\equiv x^3 +x+5 \pmod {11}$. Es sollen alle Punkte auf der Kurve berechnet werden. Es wird zuerst Gruppenordnung berechnet.
    \[
        \#E(\mathbb{F}_{11})=11+1+\sum_{x\in\mathbb{F}_{11}}\left(\frac{x^3+x+5}{\mathbb{F}_{11}}\right).
    \]
    Für $x=0,2,5,7,10$ ist $x^3+x+5$ ein Quadrat in $\mathbb F_{11}$; für die übrigen sechs Werte ist es ein Nichtquadrat. Keiner der Werte ist null. Somit gilt
    \[
        \sum_{x\in\mathbb F_{11}}
        \left(\frac{x^3+x+5}{\mathbb F_{11}}\right)=5-6=-1
    \]
    und daher $\#E(\mathbb F_{11})=12-1=11$.
    Da $11$ prim ist, erzeugt jeder Punkt außer $\mathcal O$
    die Gruppe. Wir wählen $P=(0,4)$ und berechnen seine Vielfachen.
    \[
        \begin{aligned}
            2P & = (0,4)+(0,4) = (5,5) \qquad  & 6P & = (7, 5)  \qquad   & 10P & = (0, 7) \\
            3P & = 2P+P = (10,5)        \qquad & 7P & = (2, 2) \qquad    & 11P & = \infty \\
            4P & = (2, 9)             \qquad   & 8P & = (10, 6)   \qquad & 12P & = (0,4)  \\
            5P & = (7, 6)  \qquad              & 9P & = (5, 6) \qquad    & 13P & = (5,5)  \\
        \end{aligned}
    \]

    Wie man sieht beginnt nun die Aufzählung von neuem und es handelt sich daher um eine zyklische Gruppe.

    \begin{figure}[H]
        \centering
        \includegraphics[width=0.5\linewidth]{ell_curve_finite_field.png}
        \caption{Elliptische Kurve über $\mathbb{F}_5$, eigene Darstellung}
        \label{fig:placeholder}
    \end{figure}

    Die Punkte und das Bild wurden mit einem eigenen Python Skript generiert, siehe Anhang.
\end{beispiel}
\subsection{Struktur der Punktgruppe und Torsionspunkte}
Im vorigen Beispiel hat $E(\mathbb F_{11})$ Primzahlordnung und
ist deshalb zyklisch. Für eine allgemeine elliptische Kurve über
$\mathbb F_q$ muss die Punktgruppe jedoch nicht zyklisch sein.
Um ihre mögliche Struktur zu beschreiben, betrachten wir zunächst
Punkte, deren Vielfaches der neutrale Punkt $\mathcal O$ ist.
Um das diskrete Logarithmus für ein Kryptosystem aufzubauen, müssen wir die Ordnung von einer Gruppe kennen. Dies ist deutlich herausfordernder als das Berechnen des eben gezeigten Beispiels, bei dem extra eine passende Kurve gewählt wurde. Da zuerst überprüft werden muss ob zu einem x Wert der berechnete Wert auch einen quadratischen Rest hat und damit der Punkt exestiert. Hierfür wird der Begriff der Torsionspunkte eingeführt.
\begin{definition}[$n$-Torsionspunkte]
    \label{def:n-torsionspunkte}
    Sei $E$ eine elliptische Kurve über einem Körper $K$ und
    $n\in\mathbb N$. Die Menge
    \[
        E[n]
        =
        \{P\in E(\overline K)\mid [n]P=\mathcal O\}
    \]
    heißt die Gruppe der $n$-Torsionspunkte von $E$.
\end{definition}
\begin{beispiel}
    Sei $E$ eine elliptische Kurve über $K$ mit Charakteristik ungleich 2 und 3. Es soll nun die Menge $E[2]$ bestimmt werden. Es kann die Weierstraßgleichung $y^2=x^3+Ax+B$ umgeschrieben werden als
    \[
        y^2=(x-e_1)(x-e_2)(x-e_3), \quad e_1,e_2,e_3\in\overline{K}
    \]
    Dann ist
    \[
        E[2]=\{p\in\overline{K}\mid2\cdot P=\mathcal{O}\}
    \]
    genau die Punkte, deren y-Werte verschwinden. Daher gilt dann
    \[
        E[2]=\{\mathcal{O},(e_1,0),(e_2,0),(e_3,0)\} \quad \text{und} \quad E[2]\simeq \mathbb{Z}_2\oplus \mathbb{Z}_2
    \]
\end{beispiel}

\begin{theorem}
    \label{theo:torsionsstruktur}
    Sei $E$ eine elliptische Kurve über einem Körper $K$ und sei $n\in\mathbb N$. Es ist  $\operatorname{char}(K)=p$, falls nun $p\nmid n$ oder $\operatorname{char}(K)=0$, dann gilt
    \[
        E[n]\cong \mathbb{Z}_n \oplus \mathbb{Z}_n.
    \]
    Wenn $\operatorname{char}(K)=p>0$ ist und $p\mid n$, dann schreibt man $n=p^rn'$ mit $p\nmid n'$. Dann ist
    \[
        E[n]\cong \mathbb Z_{n'} \oplus \mathbb Z_{n'} \quad\text{oder}\quad E[n]= Z_n \oplus \mathbb Z_{n'}
    \]
\end{theorem}

\begin{proof}
    Der Beweis hierzu ist sehr langwierig und findet sich daher in \cite{washington08} III Abschnitt 3.2.
\end{proof}
\begin{satz}[Hauptsatz über endliche abelsche Gruppen]
    \label{satz:hauptsatz-end-abel-gruppen}
    Sei $G$ eine endlich erzeugte abelsche Gruppe dann gilt,
    \[
        G\simeq \mathbb{Z}_{n_1}\oplus \mathbb{Z}_{n_2}\oplus\cdots\oplus\mathbb{Z}_{n_k},\quad n_1,\dots n_k\in\mathbb{Z}
    \]
    und es gilt $n_i\mid n_{i+1}$ für alle $1\geq i\geq k-1$.
\end{satz}
\begin{proof}
    Ein Beweis kann in \cite{karpf24} Satz 10.4, S.147 gefunden werden.
\end{proof}

\begin{theorem}
    \label{theo:prod-zyklischer-gruppen}
    Sei $E$ eine elliptische Kurve über einem endlichen Körper $\mathbb{F}_q$. Dann gilt,
    \[
        E(\mathbb{F}_q) \simeq \mathbb{Z}_n \text {  oder  }  E(\mathbb{F}_q) \simeq\mathbb{Z}_{n_1} \oplus \mathbb{Z}_{n_2}
    \]
    für ein $n \geq 1$ und für beliebige $n_1,n_2\geq1$ mit $n_1\mid n_2$.
    Also dann ist die Gruppe der Punkte $E(\mathbb{F}_q)$ immer eine zyklische Gruppe oder das Produkt von zwei zyklischen Gruppen.
\end{theorem}

\begin{proof}
    Nach dem Hauptsatz über endliche abelsche Gruppen \ref{satz:hauptsatz-end-abel-gruppen} ist die endliche abelsche Gruppe $E(\mathbb{F}_q)$ isomorph zu einem direkten Summe von zyklischen Gruppen
    \(E(\mathbb{F}_q)\simeq\mathbb{Z}_{n_1}\oplus \mathbb{Z}_{n_2}\oplus \dots\oplus\mathbb{Z}_{n_r}\), wobei $n_i\mid n_{i+1}$ für $i\geq 1$. Es werden nun die $n_1$-Torsionspunkte betrachtet, also alle Punkte $P$ für die gilt \(n_1P=\mathcal{O}\). Diese bilden die Gruppe
    \[
        E(\mathbb{F}_q)[n_1]
        =
        \{P\in E(\mathbb{F}_q)\mid n_1P=\mathcal{O}\}.
    \]
    Da $n_1\mid n_i$ für jedes $i$ gilt, besitzt jeder Faktor
    $\mathbb{Z}_{n_i}$ genau $n_1$ Elemente, die durch Multiplikation
    mit $n_1$ auf das neutrale Element abgebildet werden.  Da insgesamt $r$ Faktoren vorhanden sind, folgt \(\#E(\mathbb{F}_q)[n_1]=n_1^r \).
    Es ist aber \(E(\mathbb{F}_q)\subset E[n_1]\). Nach Theorem \ref{theo:torsionsstruktur} besitzt $E[n_1]$ höchstens $n_1^2$ Elemente. Somit gilt
    \[
        n_1^r
        =
        \#E(\mathbb{F}_q)[n_1]
        \leq
        \#E[n_1]
        \leq
        n_1^2.
    \]
    Für $n_1>1$ folgt daraus \( r\leq 2\) und somit kann $E(\mathbb{F}_q)$ aus höchstens zwei nichttrivialen zyklischen Faktoren bestehen. Daher gilt
    \[
        E(\mathbb{F}_q)\simeq\mathbb{Z}_n \quad\text{oder}\quad E(\mathbb{F}_q)
        \simeq
        \mathbb{Z}_{n_1}\oplus\mathbb{Z}_{n_2},
        \qquad n_1\mid n_2.
    \]
\end{proof}

\section{Frobenius-Endomorphismen und die Hasse-Schranke}
In diesem Abschnitt werden zunächst Isogenien und Endomorphismen eingeführt.Über ihre Kerne und Grade lässt sich die Anzahl der Punkte von $E(\mathbb F_q)$ mit dem Frobenius-Endomorphismus in Verbindung bringen. Dies führt zum Satz von Hasse, der die Punktzahl durch $q+1\pm2\sqrt q$ eingrenzt. Anschließend werden die Frobenius-Spur und Supersingularität betrachtet. Die Darstellung orientiert sich an Silverman \cite[S.~67, 72~f., 79, 85 und 138]{silverman09}
\cite[S.~100]{silverman06}, Washington \cite[S.~52, 98--102, 126--127 und 130--139]{washington08},Werner \cite[S.~60--73]{werner02}, sowie . Für die kryptographische Einordnung wird Güneysu, Paar und Pelzl \cite[S.~287]{guneysu24} herangezogen- Für den Beweis der Hasse Schranke wird etwas Analysis benötigt aus Pöschel\cite[S.~115]{poeschel14}.
\subsection{Isogenien}
\begin{definition}[Isogenien]
    Seien $E_1$ und $E_2$ elliptische Kurven über
    einem Körper $K$. Eine Isogenie von $E_1$ nach $E_2$
    ist ein Morphismus
    \[
        \phi:E_1\to E_2
        \qquad\text{mit}\qquad
        \phi(\mathcal O)=\mathcal O.
    \]
    Die Kurven $E_1$ und $E_2$ heißen isogen, wenn
    eine von null verschiedene Isogenie zwischen
    ihnen existiert, d.h. \(\phi(E_1)\neq\{\mathcal{O}\}\)
\end{definition}
\begin{beispiel}
    Für jedes $m\in\mathbb{Z}$ definiert man die Multiplikationsabbildung
    \[
        [m]: E\to E.
    \]
    Für $m>0$ gilt
    \[
        [m](P)=\underbrace{P+\dots+P}_{m\text{-mal}}.
    \]
    $[m]$ ein Morphismus ist. Für $m=1$ gilt
    $[1]=\operatorname{id}_E$, also ist $[1]$ ein Morphismus. Sei nun
    $[m]$ ein Morphismus. Dann gilt
    \[
        [m+1]=\mu\circ([m],\operatorname{id}_E),
    \]
    wobei $\mu:E\times E\to E$, $(P,Q)\mapsto P+Q$, nach
    \cite[III, Theorem 3.6, S.~64]{silverman09} ein Morphismus ist. Für $m<0$ gilt \[[m](P)=[-m](-P).\] Da die Negationsabbildung ebenfalls ein Morphismus ist, ist $[m]$ auch für negative $m$ ein Morphismus.
    Schließlich gilt für jedes $m\in\mathbb Z$ die Gleichung
    $[m](\mathcal O)=\mathcal O$. Daher ist $[m]:E\to E$ eine Isogenie.
\end{beispiel}
Für eine Isogenie \(\phi:E\to E\) und eine rationale Funktion $f:E_2\to K$ in $K(E_2)$ ist die Verknüpfung
\[
    \phi^*:E_1\longrightarrow K,
    \qquad f\longmapsto f\circ\phi.
\]
eine rationale Funktion in $K(E_1)$.

\begin{definition}[Grad einer Isogenie]
    Sei $\phi:E_1\to E_2$ eine Isogenie. Der Grad von $\phi$ ist der Grad der Körpererweiterung \(K(E_1)/\phi^*(K(E_2))\),
    \[
        \deg\phi=[K(E_1):\phi*(K(E_2))]
    \]
    Für die Isogenie \([0]:E_1\to E_2,P\mapsto\mathcal{O}\) setzt man $\deg[0]=0$. Eine Isogenie heißt separabel oder inseparabel, wenn ihre Körpererweiterung separabel oder inseparabel ist.
\end{definition}

\begin{lemma}
    \label{lem:deg-isogenien}
    Sei $\phi:E_1\to E_2$ eine von null verschiedene seperable Isogenie zwischen zwei elliptischen Kurven über $K$. Dann gilt
    \[
        \#\ker \phi = \deg \phi,
    \]
    wobei der Kern auf den Punkten über $\overline{K}$ betrachtet wird.
\end{lemma}

\begin{proof}
    Der Beweis hierfür lässt sich in \cite{silverman09}, S.72f.  Theorem III.4.10. (c) finden.
\end{proof}

\subsection{Frobenius-Endomorphismus}
\begin{definition}[Frobenius Abbildung]
    \label{def:frobenius}
    Sei $E$ eine elliptische Kurve, gegeben durch eine
    Weierstraßgleichung über $\mathbb F_q$, wobei
    $q=p^r$ eine Primzahlpotenz ist. Die Abbildung
    \[
        \phi_q:E(\overline{\mathbb F}_q)\to E(\overline{\mathbb F}_q),\qquad
        (x,y)\mapsto(x^q,y^q),\qquad,
    \]
    mit
    \[
        \phi_q(\mathcal O)=\mathcal O,
    \]
    heißt Frobenius-Abbildung.
\end{definition}

\begin{lemma}
    \label{lem:frobenius-fixpunkte}
    Sei $E$ eine elliptische Kurve über dem Körper $\mathbb{F}_q$. Dann gilt
    \begin{enumerate}
        \item \(\phi_q(x,y)\in E(\overline{\mathbb{F}}_q)\).
        \item \((x,y)\in E(\mathbb{F}_q)\), genau dann wenn \(\phi_q(x,y)=(x,y)\).
    \end{enumerate}
\end{lemma}
\begin{proof}
    Sei $P=(x,y)\in E(\overline{\mathbb{F}}_q)$. Damit liegt er auf der Kurve, es wird hier die allgemine Weierstraßgleichung verwendet, da der Beweis genau gleich ist.
    \[y^2+a_1xy+a_3y=x^3+a_2x^2+a_4x+a_6\] mit $a_i\in \mathbb{F}_q$.
    Nun wird die gesamte Gleichung zur $q$-ten Potenz erhoben
    \[
        (y^2+a_1xy+a_3y)^q=(x^3+a_2x^2+a_4x+a_6)^q
    \]
    Da $\operatorname{char}(\mathbb F_q)=p$ und $q=p^r$ gilt, folgt aus dem binomischen Lehrsatz $(a+b)^p=a^p+b^p$. Durch $r$-malige Anwendung erhält man $(a+b)^q=a^q+b^q$.
    \[
        (y^q)^2+a_1^q(x^qy^q)+a_3^qy^q=(x^q)^3+a_2^q(x^q)^2+a_4^qx^q+a_6^q
    \]
    Wegen Lemma \ref{lem:a^q=a} gilt $a^q=a$ für alle $a\in\mathbb{F}_q$ und daraus folgt
    \[
        (y^q)^2+a_1(x^qy^q)+a_3^qy^q=(x^q)^3+a_2^q(x^q)^2+a_4^qx^q+a_6^q
    \]
    Also enthält die Kurvengleichung auch $x^q,y^q$ und damit gilt $\phi(P)=(x^q,y^q)\in E$. Für $\mathcal{O}$ gilt nach Definition $\phi(\mathcal{O})=\mathcal{O}$ und damit gilt die Aussage direkt, somit ist die erste Aussage bewiesen.

    Sei erneut $P=(x,y)\in E(\overline{\mathbb F}_q)$, dann gilt mit Lemma~\ref{lem:a^q=a}
    \[
        x\in\mathbb F_q\iff x^q=x
        \quad\text{und}\quad
        y\in\mathbb F_q\iff y^q=y.
    \]
    Daher folgt
    \begin{align*}
        (x,y)\in E(\mathbb{F}_q) & \iff x,y\in\mathbb{F}_q         \\
                                 & \iff x^q=x  \ \text{und}\ y^q=y \\
                                 & \iff \phi_q(x,y)=(x,y)
    \end{align*}
    Für $\mathcal O$ gilt dies ebenfalls.
\end{proof}
\begin{lemma}
    \label{lem:frobenius-grad}
    Sei $E$ eine elliptische Kurve über \(\mathbb{F}_q\). Dann ist $\phi_q$ ein Endomorphismus von $E$ mit dem Grad $q$ und $\phi_q$ ist nicht seperabel.
\end{lemma}
\begin{proof}
    Die Frobenius-Abbildung $\phi_q(x,y) = (x^q, y^q)$ ist ist durch Polynome und damit insbesondere durch rationale Funktionen gegeben. Es bleibt zu zeigen, dass sie ein Gruppenhomomorphismus ist. Seien $P_1=(x_1,y_1),P_2=(x_2,y_2)\in E(\overline{\mathbb{F}}_q)$ mit $x_1\neq x_2$ und $P_1+P_2=(x_3,y_3)$. Nach den den Additionsformeln \ref{def:punktaddition} gilt
    \[
        x_3=m^2-x_1-x_2,\quad y_3=m(x_1-x_3)-y_1, \quad \text{mit } m=\frac{y_2-y_1}{x_2-x_1}.
    \]
    Erhebt man die Gleichungen zur $q$-ten Potenz, so erhält man mit Lemma~\ref{lem:freshmans-dream}
    \[
        x_3^q={m'}^2-x_1^q-x_2^q,\quad y_3^q=m'(x_1^q-x_3^q)-y_1^q, \quad \text{mit } m'=\frac{y_2^q-y_1^q}{x_2^q-x_1^q}
    \]
    Dies sind genau die Additionsformeln für die Punkte
    $\phi_q(P_1)$ und $\phi_q(P_2)$. Somit gilt
    \[
        \phi_q(P_1+P_2)
        =
        \phi_q(P_1)+\phi_q(P_2).
    \]
    Ist $P_2=-P_1$, so gilt $P_1+P_2=\mathcal O$.
    Da
    \[
        \phi_q(-P_1)=-\phi_q(P_1)
    \]
    und $\phi_q(\mathcal O)=\mathcal O$ gilt, bleibt die Homomorphieeigenschaft auch in diesem Fall erhalten. Ebenso ist der Fall, dass einer der Punkte $\mathcal O$ ist, wegen $P+\mathcal O=P$ unmittelbar erfüllt. Es bleibt die Verdopplung eines Punktes $P_1=(x_1,y_1)$ mit $y_1\neq0$. Für \(2P_1=(x_3,y_3)\) gilt
    \[
        x_3=m^2-2x_1,\quad y_3=m(x_1-x_3)-y_1\quad \text{mit } m=\frac{3x_1^2+A}{2y_1}.
    \]
    Nach erheben zur $q$-ten Potenz folgt erneut mit Lemma~\ref{lem:freshmans-dream}
    \[
        x_3^q={m'}^2-2x_1^q,\quad y_3^q=m'(x_1^q-x_3^q)-y_1^q\quad \text{mit } m'=\frac{3^q(x_1^q)^2+A^q}{2^q y_1^q}.
    \]
    Da $2,3,A\in\mathbb{F}_q$ sind, gilt mit Lemma~\ref{lem:a^q=a} $2^q=2$,$3^q=3, A^q=A$. Somit sind dies wiederum genau die Verdopplungsformeln für
    $\phi_q(P_1)$ und daher \(\phi_q(2P_1)=2\phi_q(P_1)\).  Für $y_1=0$ gilt $2P_1=\mathcal O$, sodass die Aussage ebenfalls erfüllt ist.  Damit ist $\phi_q$ ein Gruppenhomomorphismus und somit ein Endomorphismus von $E$.
    Die erste Koordinatenfunktion von $\phi_q$ ist $x^q$.
    Nach der Definition des Grades eines Endomorphismus gilt daher \(\deg(\phi_q)=q\). Schließlich gilt $q=p^r$ und $\operatorname{char}(\mathbb F_q)=p$. Daher ist
    \[
        \frac{d}{dx}x^q
        =
        qx^{q-1}
        =
        0.
    \]
    Somit ist $\phi_q$ nicht separabel.
\end{proof}
\begin{bemerkung}
    Als Konsequenz aus dem Lemma wird die Abbildung nun immer als Frobenius-Endomorphismus bezeichnet.
\end{bemerkung}
\begin{corollary}
    \label{korr:abbildung-seperabel}
    Sei $E$ eine elliptische Kurve über einem endlichen Körper $\mathbb{F}_q$ mit Charakterstik $p>0$.  Für $m,bn\in\mathbb Z$ ist der Endomorphismus
    \[
        [m]+[n]\phi_q:E\to E
    \]
    genau dann seperabel, wenn $p \nmid a$. Insbesondere ist $\operatorname{id}-\phi_q$ seperabel.
\end{corollary}

\begin{proof}
    Den Beweis findet man in \cite{silverman09}, S.79.
\end{proof}
Für $m\geq 1$ bezeichnet.
\[
    \phi_q^m
    =
    \underbrace{\phi_q\circ\cdots\circ\phi_q}_{m\text{-mal}}
\]
die $m$-fache Iteration des Frobenius-Endomorphismus. Dabei gilt
\[
    \phi_q^m(x,y)
    =
    (x^{q^m},y^{q^m}).
\]

\begin{proposition}
    \label{prop:frobenius-erweiterungskoerper}
    Sei $E$ eine elliptische Kurve über $\mathbb{F}_q$ und sei
    $m\geq 1$. Dann gelten
    \begin{enumerate}
        \item \(\ker(\phi_q^m-\operatorname{id})=E(\mathbb{F}_{q^m})\).
        \item Der Endomorphismus \(\phi^m_q-\operatorname{id}\) ist seperabel und es gilt $\#E(\mathbb{F}_{q^m})=\deg(\phi_q^m-\operatorname{id})$.
    \end{enumerate}
\end{proposition}
\begin{proof}
    \begin{enumerate}
        \item Sei $P=(x,y)\in E(\overline{\mathbb{F}_q})$. Da $\phi_q^m$ der Frobeniusendomorphismus bezüglich
              $\mathbb{F}_{q^m}$ ist, gilt mit Lemma~\ref{lem:frobenius-fixpunkte}
              \begin{align*}
                  P\in E(\mathbb{F}_{q^m}) & \iff \phi^m_q(P)=P                        \\
                                           & \iff \phi^m_q(P)=P                        \\
                                           & \iff \phi^m_q(P)-P=\mathcal{O}            \\
                                           & \iff P\in\ker(\phi^m_q-\operatorname{id})
              \end{align*}
              Daraus folgt dann \(\ker(\phi_q^m-\operatorname{id})=E(\mathbb{F}_{q^M})\).

        \item Da $\phi_q^m$ und $\operatorname{id}$ Endomorphismen von $E$ sind, ist auch \(\phi_q^m-\operatorname{id}\) ein Endomorphismus. Nach Korollar~\ref{korr:abbildung-seperabel} ist dieser separabel. Insbesondere ist er von Null verschieden und somit eine Isogenie. Für eine separable Isogenie gilt mit Lemma~\ref{lem:deg-isogenien}
              \[
                  \#\ker(\phi^m_q-\operatorname{id}) = \deg(\phi^m_q-\operatorname{id})
              \]
              und mit dem ersten Teil folgt
              \[
                  \#E(\mathbb{F}_{q^m})= \#\ker(\phi_q^m-\operatorname{id})=\deg(\phi_q^m-\operatorname{id})
              \]
    \end{enumerate}
\end{proof}

\subsection{Hasse-Schranke}
Die bisher betrachtete explizite Punktzählung ist für kleine
endliche Körper gut durchführbar, bei großen Körpern jedoch
aufwendig. Die Hasse-Schranke gibt ein Intervall für die Anzahl
der Punkte von $E(\mathbb F_q)$ an, ohne alle Punkte einzeln
bestimmen zu müssen. Das ist sowohl für die Wahl kryptographischer Kurvenparameter als auch für Schoofs Algorithmus wichtig, der die Gruppenordnung später exakt bestimmt.

Die Darstellung orientiert sich an Werner \cite[S.~60--63]{werner02}
sowie Silverman \cite[S.~138]{silverman09}
\cite[S.85,100]{silverman06}. Die kryptographische Einordnung
folgt Güneysu, Paar und Pelzl \cite[S.~287]{guneysu24}.
\begin{satz}[Satz von Hasse]
    \label{def:hasse-schranke}
    Sei $E$ eine elliptische Kurve über einem endlichen Kröper $\mathbb{F}_q$. Dann erfüllt die Ordnung von $E(\mathbb{F}_q)$,
    \[|\#E(\mathbb{F}_q)-q-1 |\leq 2\sqrt q\]
\end{satz}
Es ist nun möglich die Gruppenordnung von $E(\mathbb{F})$ rechteinfach abzuschätzen durch
\[
    -2\sqrt{q}+q+1\leq \#E(\mathbb{F}_q)\leq 2\sqrt{q}+q+1.
\]
Dieses Theorem hat nun eine weitrechende Auswirkung auf die Implementierung von elliptischen Kurven. Die Hasse-Schranke liefert damit ein Intervall für die Gruppenordnung, ohne diese exakt zu bestimmen. Insbesondere liegt \(\#E(\mathbb F_q)\) bei großen \(q\) nahe bei \(q+1\). Soll die Punktgruppe einer Kurve über einem Primkörper eine Größenordnung von \(2^{256}\) haben, muss daher auch die Primzahl \(p\) ungefähr 256 Bit groß sein (vgl. \cite[S.~287]{guneysu24}).

Die Größe $q+1-\#E(\mathbb{F}_q)$ wird auch Spur des Frobenius \ref{def:frobenius-spur} genannt.
\begin{definition}[Quadratische Form]
    \label{def:quadratische-form}
    Sei $A$ eine abelsche Gruppe. Eine Funktion
    \[
        d: A\to \mathbb{R}
    \]
    \begin{enumerate}
        \item \(d(\alpha)=d(-\alpha)\)
        \item Die Abbildung
              \[
                  A\times A \to \mathbb{R}, \quad (\alpha,\beta)\mapsto d(\alpha+\beta)-d(\alpha)-d(\beta)
              \]
              Die quadratische Form $d$ heißt positiv definit, wenn zusätzlich
        \item $d(\alpha)\geq 0 $ für alle $\alpha\in A$
        \item $d(\alpha)=0$, wenn $\alpha=0$
    \end{enumerate}
\end{definition}
\begin{corollary}
    \label{kor:gradabbildung}
    Für die elliptische Kurven $E_1$ und $E_2$ ist die Abbildung
    \[
        \deg:\operatorname{Hom}(E_1,E_2)\to\mathbb{Z}, \qquad \phi\mapsto\deg(\phi),
    \]
    eine positivie definite quadratische Form.
\end{corollary}
\begin{proof}
    Ein Beweis findet sich in \cite[Corollary III.6.3, S.~85]{silverman09}.
\end{proof}
\begin{lemma}
    \label{lem:cauchy-schwarz}
    Sei $A$ eine abelsche Gruppe und
    \[
        d:A\to \mathbb{Z}
    \]
    eine positive definite quadratische Form. Dann ist
    \[
        \mid d(\psi-\phi)-d(\phi)-d(\psi)\mid\leq 2\sqrt{d(\phi)d(\psi)}, \quad \psi,\phi\in A.
    \]
\end{lemma}
\begin{proof}
    Ein Beweis findet sich in \cite[Lemma V.1.2, S.~138]{silverman09}.
\end{proof}
Nun sind alle Vorraussetzungen gegeben, um den Beweis des Satz von Hasse durchzuführen.
\begin{proof}
    Sei $E$ eine elliptische Kurve über \(\mathbb{F}_q\). Außerdem sei \(\phi: E\to E, (x,y)\mapsto (x^q,y^q)\) der Frobenius Endomorphismus. Für ein $P=(x,y)\in E(\overline{\mathbb{F}}_q)$ gilt
    \[
        P\in E(\mathbb{F}_q) \Longleftrightarrow x,y\in \mathbb{F}_q.
    \]
    Das ist äquivalent zu
    \[
        x^q=x \quad \text{ und } \quad y^q=y.
    \]
    Also gilt
    \[
        \phi(P)=(x^q,y^q)=(x,y)=P.
    \]
    Folglich erhält man
    \[
        P\in E(\mathbb{F}_q) \Longleftrightarrow \phi(P)=P.
    \]
    Betrachtet man nun den Endomorphismus
    \(\operatorname{id}-\phi:E\to E.\) Für einen Punkt $P\in E(\overline{\mathbb{F}}_q)$ gilt
    \[
        (\operatorname{id}-\phi)(P)=\mathcal O
        \Leftrightarrow
        P-\phi(P)=\mathcal O
        \Leftrightarrow
        P=\phi(P).
    \]
    Somit folgt
    \[
        E(\mathbb{F}_q)=\ker(\operatorname{id}-\phi).
    \]
    Nach Korrolar~\ref{korr:abbildung-seperabel} ist $\operatorname{id}-\phi=[1]+[-1]\phi$ separabel, da  $p \nmid 1$ gilt. Die Abilldung ist außerdem von Null verschieden und somit eine Isogenie. Für eine separable Isogenie gilt mit Lemma~\ref{lem:deg-isogenien},
    \[
        \#E(\mathbb{F}_q)
        =\#\ker(\operatorname{id}-\phi)
        =\deg(\operatorname{id}-\phi).
    \]
    Nach Korollar~\ref{kor:gradabbildung} ist die Abbildung $\deg:End(E)\to\mathbb{Z}$ eine positiv definite quadratische Form. Setzt man in Lemma~\ref{lem:cauchy-schwarz} \( A=\operatorname{End}(E), d=\deg,\psi=\operatorname{id}\) ein, so folgt
    \[
        \left|
        \deg(\operatorname{id}-\phi)
        -\deg(\phi)
        -\deg(\operatorname{id})
        \right|
        \leq
        2\sqrt{\deg(\phi)\deg(\operatorname{id})}.
    \]

    Da
    \[
        \deg(\operatorname{id}-\phi)=\#E(\mathbb F_q),
        \qquad
        \deg(\phi)=q
        \qquad\text{und}\qquad
        \deg(\operatorname{id})=1,
    \]
    erhält man
    \[
        \left|
        \#E(\mathbb F_q)-q-1
        \right|
        \leq
        2\sqrt q.
    \]
\end{proof}

\subsection{Frobenius-Spur und Supersingularität}
Die Hasse-Schranke begrenzt die Abweichung der Gruppenordnung
von $q+1$. Diese Abweichung heißt Frobenius-Spur. Sie wird
später bei Schoofs Algorithmus zur exakten Punktzählung
bestimmt. Außerdem lässt sich mit ihr die Supersingularität
einer elliptischen Kurve charakterisieren. Diese Eigenschaft
ist für die Wahl kryptographischer Kurven relevant, da der
MOV-Angriff bei supersingulären Kurven eine Reduktion des
diskreten Logarithmusproblems auf einen endlichen Körper
kleinen Erweiterungsgrades ermöglichen kann. Die Darstellung der Frobenius-Spur und der Supersingularität orientiert sich an Werner \cite[S.~66--73]{werner02} und Washington \cite[S.~126--127, 130--139]{washington08}. Für die Einordnung des MOV-Angriffs wird Washington
\cite[S.~156--157]{washington08} herangezogen.

\begin{definition}[Supersingularität]
    Eine elliptische Kurve $E$ über $\mathbb F_q$ der
    Charakteristik $p>0$ heißt supersingulär, wenn
    \[
        E[p](\overline{\mathbb{F}}_q)=\{\mathcal{O}\}.
    \]
    Dass ist äquivlant zu der Aussage, dass eine elliptische Kurve $E$ supersingulär ist, wenn sie über dem algebraischen Abschluss keine nichttrivialen $p$ Torsionspunkte hat.
\end{definition}
Die Bezeichnung supersingulär ist von der Singularität einer algebraischen Kurve zu unterscheiden. Eine elliptische Kurve ist per Definition nichtsingulär. Supersingularität bezeichnet dagegen eine zusätzliche Eigenschaft elliptischer Kurven über Körpern positiver Charakteristik.
\begin{definition}
    \label{def:frobenius-spur}
    Sei $E$ eine elliptische Kurve über dem endlichen Körper $\mathbb{F}_q$ mit Charakteristik $p$ mit $q=p^r$. Der Wert
    \[
        t=q+1-\#E(\mathbb F_q)
    \]
    heißt Spur des Frobenius-Endomorphismus und wird mit
    \[
        t=\operatorname{Spur}(\phi_q)
    \]
    bezeichnet.
\end{definition}
\begin{proposition}[Spurkriterium]
    \label{prop:spurkriterium}
    Sei $E$ eine elliptische Kurve über dem Körper $\mathbb{F}_q$ mit $q=p^r$. Und sei $t=q+1-\#E(\mathbb{F}_q)$ erneut die Frobenius Spur. Dann ist $E$ genau dann supersingular, wenn \(t\equiv 0\mod p\). Daraus folgt \(\#E(\mathbb{F}_q)\equiv 1\mod p\).
\end{proposition}
Dies ist äquivalent zu der Aussage eine Kurve heißt supersingulär falls ihre Charakteristik $p=\operatorname{char}(\mathbb{F}_q)$ die Spur des Frobenius teilt.
\begin{proof}
    Ein Beweis findet sich in \cite[Proposition~4.31, S.~130--131]{washington08}.
\end{proof}
\begin{beispiel}
    Sei $E$ die elliptische Kurve über $\mathbb{F}_7$ mit der
    Weierstraßgleichung
    \[
        E: y^2=x^3+x.
    \]
    Die Diskriminante ist
    \[
        \Delta=-16(4\cdot 1^3)=-64\not\equiv 0\pmod 7,
    \]
    somit ist die Kurve nichtsingulär. Nun werden die Punkte nachgerechnet
    \begin{table}[H]
        \centering
        \begin{tabular}{c|c|c}
            $x$ & $x^3+x\pmod 7$ & $y$-Werte    \\
            \hline
            $0$ & $0$            & $0$          \\
            $1$ & $2$            & $\pm 3$      \\
            $2$ & $3$            & \text{keine} \\
            $3$ & $2$            & $\pm 3$      \\
            $4$ & $5$            & \text{keine} \\
            $5$ & $4$            & $\pm 2$      \\
            $6$ & $5$            & \text{keine}
        \end{tabular}
        \caption{Punkte der Kurve $E:y^2=x^3+x$ über $\mathbb{F}_7$}
        \label{tab:punkte-supersingular}
    \end{table}
    Damit gibt es sieben affine Punkte und zusätzlich den Punkt im Unendlichen. Folglich gilt \(\#E(\mathbb{F}_7)=8\) und für die Spur des Frobenius ist
    \[
        \operatorname{Spur}(\phi_q)=q+1-\#E(\mathbb{F}_q)=7+1-8=0.
    \]
    Da $7\mid 0$ gilt, ist $E$ nach dem Spurkriterium supersingulär.
\end{beispiel}
Nun wird noch ein Beispiel für den Fall Charakteristik ist 2 des Körpers angeschaut.
\begin{beispiel}
    Ein weiteres Beispiel betrifft eine elliptische Kurve der Charakterstik 2. Es ist eine Kurve \(E(\mathbb{F}_2)\) gegeben durch die Weierstraßgleichung
    \[
        y^2+y=x^3+x+1.
    \]
    Es ist für $x=0$ $y^2+y=1$ und es gibt somit keine Lösung und für $x=1$ $y^2+y=1$ und somit gibt es keine Lösung über dem Körper, also \(E(\mathbb{F}_2)=\{\mathbb{O}\}\). Es gilt für die Spur   \(\operatorname{Spur}(\phi_q)=2+1-1=2\) und ist damit durch die $\operatorname{char}(\mathbb{F}_2)=2$ teilbar.
\end{beispiel}
\begin{lemma}
    Sei $E$ eine elliptische Kurve über einem endlichen Körper $\mathbb{F}_q$. Falls $E(\mathbb{F}_q)$ supersingulär ist, so ist auch $E(\mathbb{F}_{q^k})$ supersingulär für alle $k\geq 1$.
\end{lemma}
Also die Supersingularität bleibt beim Übergang zu einem Erweiterungskörper erhalten.
\begin{proof}
    Ein Beweis findet sich bei Werner
    \cite[Lemma~3.4.2, S.~66--68]{werner02}.
\end{proof}
\begin{lemma}
    Sei $p\geq 5$ eine Primzahl. Dann ist die Ellitpische Kurve $y^2=x^3+1$ über $\mathbb{F}_p$ supersingulär, genau dann wenn $p\equiv 2 \mod 3$ und die ellitpische Kurve $y^2=x^3+x$ über $\mathbb{F}_p$ supersingulär, genau dann wenn $p\equiv 3 \mod 4$.
\end{lemma}
Das entspricht auch den Bedingungen von dem oben gewählten Beispiel.
\begin{proof}
    Ein Beweis findet sich bei Washington
    \cite[Proposition~4.37, S.~135]{washington08}.
\end{proof}

\section{Punktzählung}
Bei kleinen endlichen Körpern kann die Anzahl der Punkte einer elliptischen Kurve durch Auswertung aller möglichen $x$-Koordinaten bestimmt werden. Für große Körper ist dieses Vorgehen zu aufwendig. Um eine Kurve für kryptographische Anwendungen auszuwählen, muss ihre Gruppenordnung jedoch bekannt sein: Erst dann lässt sich prüfen, ob sie eine Untergruppe geeigneter Ordnung enthält. Dieser Abschnitt stellt zunächst Baby-Step-Giant-Step und Mestres Verfahren vor. Der Schwerpunkt liegt anschließend auf Schoofs Algorithmus. Er bestimmt die Frobenius-Spur schrittweise modulo kleiner Primzahlen und nutzt die
Hasse-Schranke, um daraus die Gruppenordnung eindeutig zu ermitteln. Die Verfahren werden hier in ihren Grundzügen dargestellt. Die Darstellung orientiert sich an Schoof \cite[S.~219--228, 233--235]{schoof95} und seiner ursprünglichen Arbeit \cite{schoof85}. Für die Erläuterung von Schoofs Verfahren wird außerdem Werner \cite[S.~63--65]{werner02} und \cite[S.~101, 123--127]{washington08} herangezogen; der verwendete chinesische Restsatz folgt Werner \cite[S.~115--116]{werner02}. Die Einordnung der Laufzeit folgt Silverman \cite[S.~46]{silverman06}.

\subsection{Baby-Step-Giant-Step Verfahren}
Für größere Primzahlen ist die direkte Bestimmung der Gruppenordnung
$\#E(\mathbb F_p)$ durch vollständiges Durchlaufen aller $x$-Werte
ineffizient. Eine effizientere Möglichkeit basiert auf der
Baby-Step-Giant-Step-Strategie von Shanks. Sei $E$ eine elliptische Kurve über $\mathbb F_p$. Nach dem Satz von Hasse gilt \( p+1-2\sqrt p \leq\#E(\mathbb F_p)\leq p+1+2\sqrt p\). Es wird zunächst ein zufälliger Punkt $P\in E(\mathbb{F}_p)$ gewählt. Anschließend wird mit Hilfe des Baby-Step-Giant-Step-Verfahrens wie in~\ref{baby-step-giant-step} eine ganze Zahl
\[
    m\in
    [p+1-2\sqrt p,\;p+1+2\sqrt p]
\]
bestimmt, für die $mP=\mathbb{O}$ gilt. Ist $m$ die einzige Zahl innerhalb des Hasse-Intervalls mit dieser Eigenschaft, so muss $m=\#E(\mathbb{F}_p)$ gelten. Exestieren dagegen zwei verschiedene Zahlen $m$ und $m'$ im Hasse-Intervall mit $mP=m'P=\mathbb{O}$, so kann die Gruppenordnung aus diesem Punkt allein nicht eindeutig bestimmt werden. In diesem Fall kann jedoch  aus $(m-m')P=\mathbb{O}$ die Ordnung von \(P\) ermittelt und das Verfahren mit einem weiteren zufälligen Punkt fortgesetzt werden. Die Ordnung von $P$ ist ein Teiler von $\#E(\mathbb{F}_p)$. Das Verfahren kann anschließend mit einem weiteren zufälligen Punkt wiederholt werden, wobei die bereits gewonnene Information über einen Teiler der Gruppenordnung die weitere Berechnung in der Regel beschleunigt.

Die Laufzeit dieses Ansatzes beträgt
\[
    O\left(p^{1/4+\varepsilon}\right)
\]
für jedes $\varepsilon>0$. Nach Schoof ist das Verfahren für
vergleichsweise größere Primzahlen praktikabel, wird jedoch bei
Primzahlen mit mehr als ungefähr 20 Dezimalstellen unpraktisch.
\subsection{Mestre's Algorithm}
Mestres Algorithmus dient dazu, die gruppentheoretischen Probleme zu
umgehen, die beim zuvor beschriebenen Baby-Step-Giant-Step-Verfahren
auftreten können. Sei dazu eine elliptische Kurve
\[
    E:y^2=x^3+Ax+B
\]
über $\mathbb F_p$ gegeben. Zusätzlich wird die quadratische Twistkurve
$E'$ von $E$ betrachtet. Für einen quadratischen Nichtrest
$g\in\mathbb F_p^\ast$ ist diese gegeben durch
\[
    E':gy^2=x^3+Ax+B.
\]
Zwischen den Gruppenordnungen der beiden Kurven besteht der Zusammenhang
\[
    \#E(\mathbb F_p)+\#E'(\mathbb F_p)=2(p+1).
\]
Kann die Gruppenordnung von $E(\mathbb F_p)$ mit dem
Baby-Step-Giant-Step-Verfahren aufgrund einer ungünstigen Gruppenstruktur
nicht eindeutig bestimmt werden, wird das Verfahren stattdessen auf die
Twistkurve $E'$ angewendet. Besitzt $E(\mathbb F_p)$ einen sehr kleinen
Gruppenexponenten, so ist dies für die Gruppe $E'(\mathbb F_p)$ im
Allgemeinen nicht gleichzeitig der Fall. Dadurch können die zuvor
beschriebenen Mehrdeutigkeiten vermieden werden.

\subsection{Schoofs Algorithm}
Für Schoofs Algorithmus werden die sogennanten Divisionspolynome benötigt, die wichtig sind für die Torsionsstruktur.
\begin{definition}[Divisionspolynome]
    Für eine elliptische Kurve
    \[
        E:y^2=x^3+Ax+B
    \]
    existieren Polynome $\psi_n$, die sogenannten
    Divisionspolynome. Für $n\geq2$ gilt für einen Punkt
    $P=(x,y)\neq\mathcal O$
    \[
        [n]P=\mathcal O
        \iff
        \psi_n(P)=0.
    \]
\end{definition}
Außerdem wird noch ein Theorem über den Frobenius benötigt.
\begin{theorem}
    Sei $E$ eine elliptische Kurve über $\mathbb{F}_q$ und $t$ die Spur des Frobenius Endomorphismus $\phi_q$. Dann gilt in
    $\operatorname{End}(E)$
    \[
        \phi_q^2-t\phi_q+q=0
    \]
\end{theorem}

\begin{proof}
    Ein Beweis findet sich in
    \cite[Theorem~4.10, S.~101]{washington08}.
\end{proof}
Sei $E$ eine elliptische Kurve über dem Körper $\mathbb{F}_q$ mit $q=p^n$ und $p$ prim und $\operatorname{char}(\mathbb{F}_q)\neq 2,3$.
Schoofs Algorithmus bestimmt die Frobenius-Spur $t=q+1-\#E(\mathbb{F}_q)$ modulo mehreren kleinen Primzahlen $l\neq p$. Seien hierzu \(l_1=2,\; l_2=3,\; l_3=5,\dots,l_r\) verschiedene kleine Primzahlen. Es seien nun $l_1=2,l_2=3,l_3=5,\dots,l_r$ die ersten $r$ Primzahlen. Nach dem Chinesischen Restsatz ~\cite[S.~115--116, Satz~6.2.1]{werner02} gilt für \(L=l_1\cdots l_r\) der Isomorphismus
\[
    \mathbb Z/L\mathbb Z
    \cong
    \mathbb Z/l_1\mathbb Z
    \times\cdots\times
    \mathbb Z/l_r\mathbb Z.
\]
Damit ist $t$ modulo $L$ durch die Restklassen
\[
    t\bmod l_1,\dots,t\bmod l_r
\]
eindeutig bestimmt. Nach dem Satz von Hasse gilt
\[
    |t|\leq 2\sqrt q
\]
und damit ist $t$ als ganze Zahl eindeutig bestimmt, sobald
\[
    L=l_1\cdots l_r>4\sqrt q
\]
gilt.

Es wird zunächst $t\mod 2$ bestimmt. Da $q$ ungerade ist, gilt
\[
    t=q+1-\#E(\mathbb F_q)
    \equiv \#E(\mathbb F_q)\pmod 2.
\]
Die Gruppenordnung $\#E(\mathbb F_q)$ ist genau dann gerade, wenn
$E(\mathbb F_q)$ einen Punkt der Ordnung $2$ enthält. Diese Punkte
haben die Form $(x,0)$. Somit ist $\#E(\mathbb F_q)$ genau dann gerade,
wenn das Polynom \(X^3+AX+B\) eine Nullstelle in $\mathbb F_q$ besitzt. Da die Nullstellen von $X^q-X$ genau die Elemente von $\mathbb F_q$ sind, besitzt $X^3+AX+B$ genau dann eine Nullstelle in $\mathbb F_q$, wenn
\[
    \operatorname{ggT}
    (X^3+AX+B,\;X^q-X)\neq 1
\]
gilt. Damit lässt sich $t\mod 2$ effizient bestimmen.
Für $l\geq 3$ ist die Bestimmung von $t\mod l$ schwieriger. Daher wird hier nur eine Idee skizziert, genauere Informationen findet man in \cite{schoof95} und \cite{schoof85}.Für ungerades $n$ ist $\psi_n$ ein Polynom in $x$. Für $P=(x,y)\in E(\overline{\mathbb{F}}_q)$ gilt
\[
    P\in E[n] \Longleftarrow \psi_n(x)=0.
\]
Der Frobenius-Endomorphismus erfüllt die Gleichung
\[
    \phi_q^2-t\phi_q+q=0 \quad P\in E(\overline{\mathbb{F}}_q)
\]
Somit gilt für jeden Punkt
$P\in E(\overline{\mathbb F}_q)$
\[
    \phi_q^2(P)-t\phi_q(P)+qP=\mathcal O.
\]
Gesucht ist nun eine Zahl $\tau\in\{0,\dots,l-1\}$ sodass die Gleichung
\[
    \phi_q^2(P)-\tau\phi_q+qp=\mathcal{O}
\]
für jeden Punkt $P$ aus der endlichen Untergruppe
\[
    E[l]=\{P\in E(\overline{\mathbb{F}}_q)\mid lp=\mathcal{O}\}
\]
gilt. Falls ein solches \(\tau\) gefunden wird, so muss jedes $P\neq\mathcal{O}$ in $E[l]$
\[
    (t-\tau)\phi_q(P)=\mathcal{O}
\]
sein. Für $P\in E[l]\setminus\{\mathcal O\}$ liegt auch
$\phi_q(P)$ in $E[l]$, dass heißt $\phi_q(P)$ hat die Ordnung $l$ in der Gruppe $E(\overline{\mathbb{F}}_q)$. Daher muss $l\mid(t-\tau)$ gelten und somit
\[
    t\equiv \tau \mod l.
\]
Es wurde die Restklasse $t\mod l$ gefunden.
Die Bestimmung von $\tau$ wird mithilfe der Gruppengesetze und der
Divisionspolynome auf Polynomgleichungen zurückgeführt. Für
$\tau=0,\dots,l-1$ wird überprüft, für welchen Wert die entsprechende
Gleichung auf den $l$-Torsionspunkten erfüllt ist. Dieser Wert liefert
die gesuchte Restklasse $t\bmod l$.

Schoofs Algorithmus ist deterministisch und besitzt eine Laufzeit,
die polynomial in $\log q$ ist. Damit kann die Gruppenordnung
$\#E(\mathbb F_q)$ effizient bestimmt werden.
In der Praxis ist der ursprüngliche Schoof-Algorithmus jedoch aufgrund
der hohen Grade der Divisionspolynome vergleichsweise aufwendig.
Atkin und Elkies entwickelten Verbesserungen, die zur sogenannten
SEA-Methode führen und die Berechnungen deutlich effizienter machen.
Die Gruppenordnung lässt sich effizient bestimmen, ohne dass damit das diskrete Logarithmusproblem auf der Kurve effizient lösbar wird.

